% ID: FIS-LEN-CEN-001
% Titolo: Oggetto lanciato verso l'alto da una quota iniziale
% Disciplina: Fisica
% Area: Lavoro ed energia
% Argomento: Conservazione energia
% Sottoargomento: Energia meccanica nel campo gravitazionale
% Classe: Terza liceo scientifico
% Difficolta: 3
% Tipo: problema_numerico
% Risultato: h_max circa 4,8 m; v circa 5,9 m/s
% Tempo_stimato: 8 min
% Competenze: conservazione dell'energia meccanica, interpretazione della quota, calcolo algebrico
% Prerequisiti: energia cinetica, energia potenziale gravitazionale, velocita
% Tag: fisica, lavoro_energia, conservazione_energia, energia_meccanica, gravita
% Fonte: esempio_originale
% Autore: OpenAI Codex
% Licenza: CC-BY-SA-4.0

\begin{esercizio}
Un oggetto viene lanciato verticalmente verso l'alto da una quota iniziale
\(h_0=1{,}5\,\mathrm{m}\) rispetto al suolo, con velocita iniziale
\(v_0=8{,}0\,\mathrm{m/s}\). Trascuriamo la resistenza dell'aria.

Determina la quota massima raggiunta dall'oggetto e la sua velocita quando
si trova alla quota \(h=3{,}0\,\mathrm{m}\). Usa
\(g=9{,}8\,\mathrm{m/s^2}\).
\end{esercizio}

\begin{soluzione}
L'energia meccanica si conserva:
\[
\frac{1}{2}mv^2+mgh=\text{costante}.
\]

Alla quota massima la velocita e nulla. Confrontiamo lo stato iniziale con
quello alla quota massima:
\[
\frac{1}{2}mv_0^2+mgh_0=mgh_\text{max}.
\]

La massa si semplifica:
\[
h_\text{max}=h_0+\frac{v_0^2}{2g}.
\]

Sostituendo i dati:
\[
h_\text{max}=1{,}5+\frac{(8{,}0)^2}{2\cdot 9{,}8}
\simeq 1{,}5+3{,}3
=4{,}8\,\mathrm{m}.
\]

Alla quota \(3{,}0\,\mathrm{m}\) vale ancora:
\[
\frac{1}{2}mv_0^2+mgh_0=\frac{1}{2}mv^2+mgh.
\]

Quindi:
\[
v^2=v_0^2+2g(h_0-h).
\]

\[
v^2=(8{,}0)^2+2\cdot 9{,}8\cdot (1{,}5-3{,}0)
=64-29{,}4=34{,}6.
\]

\[
v=\sqrt{34{,}6}\simeq 5{,}9\,\mathrm{m/s}.
\]

Il valore trovato e il modulo della velocita; alla quota indicata l'oggetto
puo avere quel modulo sia in salita sia in discesa.
\end{soluzione}
